% ============================================================
%  Part IV: 反重力と重力工学
%  wb_textbook_v2_antigravity.tex
%  Chapters 11--13
% ============================================================

% ============================================================
% CHAPTER 11
% ============================================================
\chapter{$\Geff = -G$ 定理と反重力の物理}
\label{ch:antigravity}

\begin{keyresult}[本章の主題]
Berry位相 $\phi = \pi$ の物質内部では、有効重力定数が符号反転する：$\Geff = -G$。
ただし、この反重力は\textbf{局所限定}であり、宇宙論的de Sitter解は生じない。
\end{keyresult}

% ------------------------------------------------------------
\section{中心定理の復習}
\label{sec:11.1}

Part III で導出した鍵となる論理連鎖を振り返ろう。

\begin{enumerate}[label=(\arabic*)]
  \item 物質中のバンド構造が\textbf{$\pi$-Berry位相}を生む（トポロジカル絶縁体等）。
  \item Berry位相 $\phi$ がスペクトル三つ組の自己同型 $f_k \to e^{-i\phi(d-k)/2}f_k$ を誘導する。
  \item スペクトル作用を通じて、分配関数が $K(t) \to -\zeta(t)$ と符号反転する。
  \item 結果として有効ニュートン定数が $\Geff = G\cos\phi$ となる。
  \item $\phi = \pi$ のとき $\cos\pi = -1$ より $\Geff = -G$。
\end{enumerate}

\begin{hsbox}
「Berry位相」とは、量子力学の状態がパラメータ空間を一周したときに蓄積する幾何学的位相のこと。
通常の波動関数の位相 $e^{i\theta}$ とは異なり、Berry位相はバンド構造の\textbf{トポロジー}（幾何学的形状）で決まる。
$0$ か $\pi$ かの二値しか取れない場合が多く、$\pi$ の場合を「トポロジカル」と呼ぶ。
\end{hsbox}

% ------------------------------------------------------------
\section{新しい理解：$\cos(2\phi)$ と $\cos(\phi)$ の非対称性}
\label{sec:11.2}

ここで、Part III の結果をさらに掘り下げる。Berry位相 $\phi$ に対して、
宇宙定数 $\Lambda$ と重力定数 $G$ は\textbf{異なる位相依存性}を持つ。

\begin{dfnbox}[$\Lambda$ と $G$ の位相依存性]
\begin{align}
  \Lambda_{\mathrm{eff}}(\phi) &\propto \cos(2\phi), \label{eq:lambda-phi} \\
  G_{\mathrm{eff}}(\phi) &= G\cos(\phi). \label{eq:g-phi}
\end{align}
\end{dfnbox}

この違いは本質的である。$\phi = \pi$ を代入すると：
\begin{align}
  \Lambda_{\mathrm{eff}}(\pi) &= \Lambda \cdot \cos(2\pi) = \Lambda \cdot (+1) = +\Lambda, \\
  G_{\mathrm{eff}}(\pi) &= G \cdot \cos(\pi) = G \cdot (-1) = -G.
\end{align}

\begin{keyresult}[$\phi = \pi$ での非対称性]
Berry位相 $\phi = \pi$ では：
\begin{itemize}
  \item $G$ は\textbf{反転する}（$\Geff = -G$）：重力が斥力に変わる。
  \item $\Lambda$ は\textbf{反転しない}（$\Lambda_{\mathrm{eff}} = +\Lambda$）：宇宙定数は不変。
\end{itemize}
この非対称性が $\cos$ の引数の違い（$\phi$ vs $2\phi$）から生じることに注意せよ。
\end{keyresult}

\begin{hsbox}
三角関数の復習：$\cos(\pi) = -1$ だが $\cos(2\pi) = +1$。
角度を $\pi$（半回転）動かすと $\cos$ は符号が反転するが、
角度を $2\pi$（全回転）動かすと元に戻る。
この単純な数学的事実が「$G$ だけ反転し $\Lambda$ は変わらない」という深い物理を生む。
\end{hsbox}

\subsection{$\phi = \pi$ に de Sitter 解なし}

$\Lambda$ が反転しないという事実は重要な帰結を持つ。
通常、正の宇宙定数 $\Lambda > 0$ は de Sitter 時空（指数関数的膨張解）を与える。
もし $\phi = \pi$ で $\Lambda$ が負に反転すれば、反 de Sitter 時空が生じ、
宇宙全体の構造が変わってしまう。

しかし式 \eqref{eq:lambda-phi} により $\Lambda_{\mathrm{eff}}(\pi) = +\Lambda$ のまま不変なので、
$\phi = \pi$ 領域には通常の de Sitter 解は残る。ただし $G < 0$ なので
Einstein 方程式の性質が根本的に変わり、\textbf{通常の de Sitter 解は成立しない}。

\begin{warningbox}[局所限定の原理]
$\phi = \pi$ の反重力領域は\textbf{物質内部のみ}に限定される。
真空中では $\phi$ は安定点 $\phi = 0$ に緩和するため、
マクロスコピックな反重力領域を維持するには Berry 位相 $\pi$ を支える物質が不可欠。
\end{warningbox}

% ------------------------------------------------------------
\section{物質Berry位相 = 重力スカラー場への境界条件}
\label{sec:11.3}

なぜ物質が $\phi = \pi$ を「固定」できるのかを、2つのアナロジーで理解しよう。

\subsection{強磁性体のアナロジー}

強磁性体（鉄やコバルト）では、磁化ベクトル $\mathbf{M}$ は自由に回転できるわけではない。
結晶構造が「容易軸」を定め、$\mathbf{M}$ はその方向に固定される
（\textbf{磁気異方性}）。外部磁場がなくても、結晶が磁化の向きを決める。

同様に、重力スカラー場 $\phi$ は真空中では自由に振る舞うが、
物質内部では\textbf{バンド構造}が $\phi$ の値を固定する。
トポロジカル物質では、バンドのトポロジーが $\phi = \pi$ を強制する。

\begin{center}
\begin{tabular}{lll}
\toprule
& 磁性体 & Berry位相物質 \\
\midrule
場 & 磁化 $\mathbf{M}$ & スカラー場 $\phi$ \\
固定するもの & 結晶の対称性 & バンドのトポロジー \\
固定される値 & 容易軸方向 & $\phi = 0$ or $\pi$ \\
\bottomrule
\end{tabular}
\end{center}

\subsection{ヒッグス場のアナロジー}

超伝導体内部では、ヒッグス場（クーパー対の凝縮体）$\phi_H$ が
真空とは異なる値を取る。これによりマイスナー効果が生じ、磁場が排除される。

同様に、$\pi$-Berry物質の内部では $\phi = \pi$ が固定され、
真空の $\phi = 0$ とは異なる「重力的超伝導状態」が実現する。

\begin{dfnbox}[境界条件としてのBerry位相]
物質のBerry位相 $\phi_B$ は、重力スカラー場 $\phi$ に対する\textbf{境界条件}として機能する：
\[
  \phi\big|_{\text{物質内部}} = \phi_B.
\]
トポロジカル物質では $\phi_B = \pi$（量子化されているので中間値を取れない）。
\end{dfnbox}

% ------------------------------------------------------------
\section{Nielsen-Ninomiya定理のボソン系抜け穴}
\label{sec:11.4}

ここで決定的に重要な問題がある。$\pi$-Berry位相を持つ物質は\textbf{実在するのか}？

\subsection{フェルミオン系の制約}

Nielsen-Ninomiya定理（1981年）は、格子上のフェルミオン系に厳しい制約を課す。
カイラル対称性を保ったまま格子正則化を行うと、\textbf{ダブラー}と呼ばれる
余分な自由度が必ず現れる。

この定理の帰結として、フェルミオン（電子）のバンド構造では、
個々のバンドが $\pi$-Berry位相を持つことはあっても、
\textbf{バルク全体}の Berry 位相を $\pi$ にすることは困難である。
Weyl ノードは常にペアで現れ、全体としての Berry 位相は $0$ に帰着しやすい。

\subsection{ボソン系の抜け穴}

しかし Nielsen-Ninomiya 定理は\textbf{フェルミオン系}に対する制約であり、
\textbf{ボソン系}には適用されない。フォノン（格子振動）やマグノン（スピン波）は
ボソンであり、この制約から自由である。

\begin{keyresult}[ボソン系のBerry位相 $\pi$]
フォノニック結晶（音響メタマテリアル）では：
\begin{itemize}
  \item Nielsen-Ninomiya 定理の制約を受けない。
  \item バルク全体で Berry 位相 $\pi$ を実現可能。
  \item Wilson loop 計算で $\pi$ が\textbf{確認済み}（2球体構造、反転対称性の破れ）。
\end{itemize}
\end{keyresult}

\begin{hsbox}
フェルミオンとボソンの違い：電子（フェルミオン）は同じ場所に2個までしか入れない（パウリの排他原理）。
一方、フォノン（音の振動、ボソン）にはそのような制約がない。
Nielsen-Ninomiya 定理はパウリの排他原理と深く関係するので、ボソンには効かない。
\end{hsbox}

これがPhase P実験（第14章）でフォノニック結晶を使う根本的理由である。


% ============================================================
% CHAPTER 12
% ============================================================
\chapter{重力ダイヤル $\Geff = G\cos(\phi)$}
\label{ch:gravity-dial}

\begin{keyresult}[本章の主題]
Berry位相 $\phi$ を $0$ から $\pi$ まで連続的に制御できれば、
重力定数を $+G$ から $-G$ まで自在に調整する「重力ダイヤル」が実現する。
\end{keyresult}

% ------------------------------------------------------------
\section{Berry位相の連続制御}
\label{sec:12.1}

前章では $\phi = 0$（通常重力）と $\phi = \pi$（反重力）の二値を議論した。
しかし原理的には Berry 位相は\textbf{連続的に}制御できる。

連続制御が可能な系の条件：
\begin{enumerate}[label=(\roman*)]
  \item バンドギャップが閉じないこと（トポロジカル相転移を避ける）。
  \item 外部パラメータ（磁場、ゲート電圧、歪み等）で $\phi$ を掃引できること。
  \item 量子化条件が破れていること（対称性が十分に低いこと）。
\end{enumerate}

\begin{hsbox}
「ダイヤル」という比喩を使っているのは、ラジオのダイヤルのように
連続的に値を変えられることを強調するため。$\phi = 0$（重力ON）から
$\phi = \pi$（反重力ON）まで、好きな値に合わせられる。
\end{hsbox}

% ------------------------------------------------------------
\section{$\Geff(\phi)$ テーブル}
\label{sec:12.2}

$\Geff = G\cos\phi$ の代表的な値を表にまとめる。

\begin{center}
\begin{tabular}{ccrl}
\toprule
$\phi$ & $\cos\phi$ & $\Geff/G$ & 物理的意味 \\
\midrule
$0$      & $+1$   & $+1$    & 通常重力（標準物理） \\
$\pi/6$  & $+\frac{\sqrt{3}}{2}$ & $+0.866$ & わずかに軽減 \\
$\pi/3$  & $+\frac{1}{2}$ & $+0.5$  & 重力半減 \\
$\pi/2$  & $0$    & $0$     & \textbf{無重力}（重力遮断） \\
$2\pi/3$ & $-\frac{1}{2}$ & $-0.5$  & 反重力半分 \\
$5\pi/6$ & $-\frac{\sqrt{3}}{2}$ & $-0.866$ & 強い反重力 \\
$\pi$    & $-1$   & $-1$    & \textbf{完全反重力} \\
\bottomrule
\end{tabular}
\end{center}

$\phi = \pi/2$ の「無重力」は特に興味深い。
重力が完全にゼロになる領域は、宇宙ステーション内の微小重力とは本質的に異なる。
後者は自由落下による見かけの無重力だが、前者は$G$自体がゼロになる真の無重力である。

% ------------------------------------------------------------
\section{4つの装置コンセプト}
\label{sec:12.3}

Berry位相 $\phi$ を制御する4種類の実験装置を概説する。

\subsection{$\pi$-接合プレート（パッシブ型、推定\$40k）}

最もシンプルなアプローチ。$\phi = \pi$ を持つトポロジカル絶縁体
（例：Bi$_2$Se$_3$）を薄板状に加工し、カシミール効果実験の一方の板として使用する。

\begin{itemize}
  \item 制御パラメータ：なし（$\phi = \pi$ 固定）。
  \item 利点：外部電源不要、単結晶があれば即実験可能。
  \item 欠点：$\phi = 0$ か $\pi$ の二値のみ、連続制御不可。
  \item 推定コスト：試料 + カシミール測定装置で約 \$40,000。
\end{itemize}

\subsection{FQHE重力ダイヤル（連続制御）}

分数量子ホール効果（FQHE）系では、充填率 $\nu$ を磁場で制御することにより、
準粒子のBerry位相を連続的に変化させられる。

\begin{itemize}
  \item 制御パラメータ：外部磁場 $B$（充填率 $\nu \propto 1/B$）。
  \item 利点：$\phi$ の連続制御が原理的に可能。
  \item 欠点：極低温（$\sim$ mK）、強磁場（$\sim$ 10\,T）が必要。
  \item 研究段階：概念設計のみ。
\end{itemize}

\subsection{グラフェン重力ダイヤル（ゲート電圧制御）}

二層グラフェンでは、層間のゲート電圧によりバレーBerry位相を制御できる。

\begin{itemize}
  \item 制御パラメータ：ゲート電圧 $V_g$。
  \item 利点：室温動作の可能性、スケーラブル。
  \item 欠点：バレー自由度の Berry 位相が重力スカラー場 $\phi$ と直接結合するか未確認。
  \item 研究段階：理論検討中。
\end{itemize}

\subsection{歪みトポロジカル絶縁体（最安）}

機械的歪みによりバンド構造を変形し、トポロジカル相転移を制御する。

\begin{itemize}
  \item 制御パラメータ：機械的歪み $\epsilon$。
  \item 利点：特殊な装置不要、最も安価。
  \item 欠点：歪みの精密制御が困難、ヒステリシスの可能性。
  \item 研究段階：候補材料の選定中。
\end{itemize}

% ------------------------------------------------------------
\section{真空量子化との関係}
\label{sec:12.4}

上記の「連続ダイヤル」には重要な制約がある。

真空中の $\phi$ は\textbf{量子化}されている。有効ポテンシャル $V(\phi)$ が
$\phi = 0$ と $\phi = \pi$ に極小を持ち、中間値は不安定である。

\begin{warningbox}[真空量子化の制約]
平衡状態では、Berry位相は $\phi = 0$ または $\phi = \pi$ の\textbf{二値のみ}が安定。
中間値 $0 < \phi < \pi$ は過渡的（準安定）であり、
持続的に維持するには外部エネルギーの投入が必要。
\end{warningbox}

これは強磁性体の磁化 $M$ が「上」か「下」のどちらかを向くのと類似している。
中間の角度は外部磁場で一時的に実現できるが、磁場を切ると上か下に戻る。

つまり、重力ダイヤルの中間値（$\phi = \pi/2$ の無重力状態など）は
外部制御下でのみ実現可能であり、エネルギーコストを伴う。
パッシブに維持できるのは $\phi = 0$（通常重力）と $\phi = \pi$（完全反重力）のみである。


% ============================================================
% CHAPTER 13
% ============================================================
\chapter{ドメインウォールとBerry位相インフレーション}
\label{ch:domain-walls}

\begin{keyresult}[本章の主題]
$\phi = 0$ と $\phi = \pi$ の境界には「重力ドメインウォール」が形成される。
さらに、初期宇宙での $\phi = \pi \to 0$ の相転移がインフレーションを駆動し、
ダークエネルギーと同一のポテンシャルから説明できる可能性がある。
\end{keyresult}

% ------------------------------------------------------------
\section{ドメインウォール}
\label{sec:13.1}

$\phi = 0$（通常重力）と $\phi = \pi$（反重力）は、有効ポテンシャル $V(\phi)$ の
2つの極小に対応する。この2つの真空の間には\textbf{ドメインウォール}（壁）が存在する。

\subsection{キンク解}

1次元のスカラー場理論において、2つの極小を結ぶ静的解は
\textbf{キンク}（kink）と呼ばれる。$\phi$ のドメインウォールも同型の解を持つ：
\[
  \phi(x) = \frac{\pi}{2}\left[1 + \tanh\left(\frac{x}{\ell_w}\right)\right],
\]
ここで $\ell_w$ はウォールの幅である。$x \to -\infty$ で $\phi \to 0$、
$x \to +\infty$ で $\phi \to \pi$ となる。

\subsection{ウォールの幅}

ウォールの幅はポテンシャルのバリア高さ $\delta$ で決まる：
\[
  \ell_w \sim \frac{1}{\sqrt{2\delta}}.
\]

\begin{keyresult}[ドメインウォールの幅]
$\delta$ がダークエネルギースケール（$\sim \Lambda_{\mathrm{obs}}$）の場合：
\[
  \ell_w \sim \frac{1}{H_0} \sim \text{ハッブル半径} \sim 10^{26}\,\text{m}.
\]
つまり、重力ドメインウォールの幅は\textbf{観測可能な宇宙のサイズ}程度になる！
\end{keyresult}

\begin{hsbox}
ドメインウォールとは、磁石の「磁区」の境界と同じ概念。
鉄の磁石の中では、磁化の向きが違う領域（N極向きとS極向き）が隣り合っており、
その境界では磁化が徐々に回転する。この境界がドメインウォール。
同様に、$\phi = 0$ の領域と $\phi = \pi$ の領域の境界で $\phi$ が徐々に変化する。
\end{hsbox}

\subsection{重力ドメインウォール}

このウォールを横切ると、$G$ の符号が反転する。
ウォールの片側では通常の引力重力、反対側では斥力重力が支配する。

ウォールの面エネルギー密度は
\[
  \sigma \sim \sqrt{\delta}\, M_{\mathrm{Pl}}^2
\]
であり、$\delta \sim \Lambda_{\mathrm{obs}}$ の場合は非常に小さい。

% ------------------------------------------------------------
\section{Berry位相インフレーション}
\label{sec:13.2}

ドメインウォールの物理は、初期宇宙の\textbf{インフレーション}と
現在の\textbf{ダークエネルギー}を統一的に説明する可能性を開く。

\subsection{基本的アイデア}

\begin{enumerate}[label=(\arabic*)]
  \item \textbf{初期宇宙}：$\phi \approx \pi$（反重力相）に閉じ込められている。
  \item \textbf{相転移}：量子トンネリングにより $\phi = \pi \to \phi = 0$ へ遷移。
  \item \textbf{現在}：$\phi = 0$ の底にいるが、残余のポテンシャルエネルギー $V(0) = \delta$ がダークエネルギーとして観測される。
\end{enumerate}

\subsection{Einstein フレームでのバリア}

Jordan フレームでのポテンシャル $V(\phi)$ を Einstein フレームに変換すると、
$\phi = \pi$ 近傍でバリアが\textbf{発散}する。これは共形変換
$\tilde{g}_{\mu\nu} = |\cos\phi|\, g_{\mu\nu}$ において $\cos\pi = -1 \to 0$ の
極を横切るためである。

この発散バリアのため、古典的な $\phi$ のロールオーバーは不可能であり、
遷移は\textbf{Coleman-De Luccia トンネリング}によってのみ起こる。

\begin{dfnbox}[Coleman-De Luccia トンネリング]
偽真空（$\phi \approx \pi$）から真真空（$\phi = 0$）への
量子力学的トンネリング過程。泡（バブル）核生成として記述され、
バブルが膨張して偽真空を侵食する。重力効果を含めた計算では、
バリアが高いほどトンネリング率が指数関数的に抑制される。
\end{dfnbox}

\subsection{インフレーション + ダークエネルギーの統一}

Berry位相インフレーションの最も魅力的な特徴は、\textbf{同じポテンシャル $V(\phi)$} から
インフレーションとダークエネルギーの両方が説明できることである。

\begin{keyresult}[同一ポテンシャルからの二重説明]
\begin{itemize}
  \item \textbf{インフレーション期}：$\phi \approx \pi$ の偽真空のエネルギー $V(\pi)$ が
        指数関数的膨張を駆動。
  \item \textbf{現在のダークエネルギー}：$\phi = 0$ の真空エネルギー $V(0) = \delta$ が
        加速膨張を駆動。
  \item 両者は $V(\phi)$ の異なる極小の値に対応。
\end{itemize}
\end{keyresult}

\subsection{Slow-roll パラメータと観測予測}

$\delta$ がダークエネルギースケールで決まるため、インフレーション期の
slow-roll パラメータも $\delta$ に依存する：
\begin{align}
  \epsilon &= \frac{M_{\mathrm{Pl}}^2}{2}\left(\frac{V'}{V}\right)^2, \\
  \eta &= M_{\mathrm{Pl}}^2 \frac{V''}{V}.
\end{align}

これらからスカラースペクトル指数 $n_s$ とテンソル-スカラー比 $r$ が予測される：
\begin{align}
  n_s &= 1 - 6\epsilon + 2\eta, \\
  r &= 16\epsilon.
\end{align}

\begin{warningbox}[$\delta$ 依存性に関する注意]
$n_s$ と $r$ の具体的な数値は $\delta$ の正確な値に強く依存する。
$\delta$ はダークエネルギー密度 $\rho_\Lambda \sim 10^{-47}\,\mathrm{GeV}^4$ と
関係するが、$V(\phi)$ の全体的な形状（kinetic term の正規化を含む）が
まだ厳密に導出されていないため、定量的予測は現時点では不確実である。
これは第18章で未解決問題として議論する。
\end{warningbox}

\begin{hsbox}
インフレーションとは、宇宙が誕生直後（$\sim 10^{-36}$ 秒）に
指数関数的に膨張した時期のこと。これにより、宇宙がなぜこんなに一様で平坦なのかが説明できる。
ダークエネルギーとは、現在の宇宙の加速膨張を引き起こしている未知のエネルギー。
この2つが同じ原因（同じポテンシャル $V(\phi)$）で説明できるなら、
物理学の2大謎が一挙に解決する可能性がある。
\end{hsbox}
